\documentclass[11pt,a4paper]{article}

\usepackage[utf8]{inputenc}

\usepackage[T1]{fontenc}

\usepackage[polish]{babel}

\usepackage{lmodern}

\usepackage{geometry}

\geometry{margin=27mm}

\usepackage[hidelinks]{hyperref}

\title{Walidacja miary postrzeganej sprawczości humanoida}

\author{Mateusz Pomianek}

\date{Szkic naukowy, wersja 1.0.0 | 28 listopada 2026}

\begin{document}

\maketitle

\begin{abstract}

Proponowana miara główna to różnica wyniku logarytmicznego prognoz odpowiedzi dla konkurencyjnych struktur czynnikowych.

\end{abstract}

\section{Wprowadzenie}

Sprawczość i zdolność do doświadczania dają się rozdzielić w badaniach percepcji umysłu (Gray et al., 2007). Godspeed dostarcza miar sąsiednich, w tym antropomorfizacji i inteligencji (Bartneck et al., 2009).

\section{Cel i pytania badawcze}

Hipotezy robocze wymagają zatwierdzenia przed rozpoczęciem właściwej weryfikacji.

\paragraph{H1.} W niezależnej próbie model rozdzielający sprawczość i doświadczanie będzie lepiej prognozował odpowiedzi na wspólną pulę pozycji niż model jednoczynnikowy.

\paragraph{H2.} Latentna sprawczość będzie silniej związana z oceną zdolności robota do osiągnięcia celu niż z oceną sympatii.

\section{Materiał i metody}

Pierwszorzędowy wskaźnik: różnica wyniku logarytmicznego prognoz odpowiedzi dla konkurencyjnych struktur czynnikowych.

W próbie konstrukcyjnej należy przeprowadzić analizę treści i eksplorację, a w próbie potwierdzającej ordinalną CFA z zamrożoną specyfikacją. Porównanie predykcyjne na wstrzymanych osobach obejmie ten sam zestaw pozycji. Przed porównaniem form prezentacji sprawdzi się inwariancję.

Przed badaniem z udziałem ludzi należy zatwierdzić wymagania etyczne, zasady zgody i ochronę danych.

\section{Plan analizy}

Zamrożenie estymandu, reguł oceny jakości i obsługi braków nastąpi przed dostępem do danych wynikowych. Raport obejmie oszacowanie kontrastu głównego z przedziałem niepewności i pełne mianowniki. Jeśli przewidziano hipotezy drugorzędowe, przed analizą należy ustalić sposób kontroli wielokrotności.

\section{Wyniki}

W niniejszym pliku nie umieszczono wyników empirycznych, statystyk ani nieprzeprowadzonych analiz. Sekcja wymaga uzupełnienia po uzyskaniu i weryfikacji materiału dowodowego.

\section{Dyskusja i ograniczenia}

Ograniczenia obejmują możliwe błędy pomiaru, dobór próby, alternatywne mechanizmy i ograniczoną możliwość uogólniania. Korelacja nie dowodzi przyczynowości; poprawność techniczna nie jest dowodem znaczenia klinicznego.

\section{Wnioski}

Hipotezy i sposób pomiaru stanowią podstawę oceny wykonalności oraz przygotowania właściwego protokołu.

\section*{Bibliografia}

{\setlength{\parindent}{0pt}
\setlength{\parskip}{5pt}
\hangindent=1.27cm
Bartneck, C., Kulić, D., Croft, E., \& Zoghbi, S. (2009). Measurement instruments for the anthropomorphism, animacy, likeability, perceived intelligence, and perceived safety of robots. \textit{International Journal of Social Robotics, 1}(1), 71--81. \url{https://doi.org/10.1007/s12369-008-0001-3}\par
\hangindent=1.27cm
Gray, H. M., Gray, K., \& Wegner, D. M. (2007). Dimensions of mind perception. \textit{Science, 315}(5812), 619. \url{https://doi.org/10.1126/science.1134475}\par
}

\end{document}
